% SPDX-FileCopyrightText: Copyright (c) 2026 NVIDIA CORPORATION & AFFILIATES. All rights reserved.
% SPDX-License-Identifier: Apache-2.0
%
% ROLE: define the decision (which worker serves a request), what the router can observe, how load
%   is applied, and the objective every policy is scored by. Everything later refers back here.
% CLAIMS: definitional, plus the reasons for each definition (slowdown rather than raw latency,
%   windowed rather than makespan goodput, clipped log-ratio rather than mean ratio).
% EVIDENCE: CR/CONTRACT.md (Deployment, A2, A5.2, A14, A17.2); CR/config/engine.json;
%   WT/benchmarks/learned_routing/learned_routing/{goodput.py,e0.py,train.py};
%   CR/facts/{calibration.json,setup.json,build_workloads.json}; CR/literature/LESSONS.md LR-01,
%   LR-08, LR-09; WT/lib/router-plugins/builtin/src/learned_choice/FEATURES.md.
% STATUS: definitions settled (frozen at calibration, lr-cells-v5).

\section{Problem and objective}
\label{sec:setup}

This section fixes what a routing policy decides, what it may observe, how load is applied, and the
score by which every policy, learned or heuristic, is compared.

\subsection{The routing decision}
\label{sec:decision}

\paragraph{Deployment.}
% src: CR/CONTRACT.md "Deployment (operator-fixed)"; CR/config/engine.json mock_engine_args;
%      CR/facts/setup.json kv_capacity
A deployment runs $\Nw$ identical aggregated workers, each a vLLM engine serving Qwen3-32B with
tensor parallelism 2 on H100 SXM GPUs, so prefill and decode share a worker. Each worker holds
\num{18863} KV blocks of 16 tokens, \num{301808} tokens in all. The study covers
$\Nw \in \{2, 4, 6, 8, 16, 32\}$; heterogeneous worker sets are out of scope (Amendment A5.2).
The router never holds a request back: \code{router\_queue\_threshold} is left unset, which
disables router-side queueing; an ablation that tuned it jointly found no measurable gain
(\cref{sec:res-ladder}).
% src: WT/lib/kv-router/src/scheduling/config.rs test queueing_enabled_reflects_synthetic_threshold
%      ("With default config, queueing is disabled"); binding default router_queue_threshold=None
%      (WT/lib/bindings/python/rust/llm/entrypoint.rs:134)

\paragraph{A sequential discrete choice.}
When request $j$ arrives at time $a_j$, the router's worker-selection policy $\pol$ picks one
worker from the eligible candidate set $\cand \subseteq \{1,\dots,\Nw\}$. Three properties shape
the problem.
\begin{itemize}
  \item \emph{The choice is sequential.} The chosen worker computes the uncached part of $j$'s
    prompt and keeps the resulting KV blocks, which later requests that share the prefix can reuse
    and which may evict other prefixes' blocks. Request $j$ also adds to that worker's prefill and
    decode load until it completes. A decision's cost or benefit therefore lands on other requests,
    often much later in simulated time, and no per-decision label exists. Policies are scored on
    whole episodes (\cref{sec:objective}).
  \item \emph{The worker set varies.} $\Nw$ differs across deployments, and $\cand$ can be smaller
    than $\Nw$ when the host marks workers ineligible. A policy defined on a fixed-length vector of
    all workers' states would need retraining for every $\Nw$. A policy that scores each candidate
    with one shared function and picks the best applies to any $|\cand|$. This is the
    discrete-choice form of the learned policy (\cref{sec:model}), and it is what lets the test
    set hold out worker counts that training never saw.
  \item \emph{The choice is made online.} It runs once per request on the router's request path,
    from information available at $a_j$, so it must be cheap and must not depend on the future.
\end{itemize}

\paragraph{What the router observes.}
% src: WT/lib/router-plugins/builtin/src/learned_choice/FEATURES.md "Feature set v1",
%      "Excluded inputs", "Information parity with the live router"
At decision time a policy sees, through the router plugin API (\cref{sec:api-components}):
\begin{itemize}
  \item the request: its prompt length $\ISL_j$, its prefix block hashes and an optional session
    identifier;
  \item per candidate, the router's cache estimate: blocks of overlap with the request's prefix and
    the estimated number of cached tokens;
  \item per candidate, load the router tracks from its own dispatches: active prefill tokens,
    decode blocks and active requests;
  \item per candidate, its configured KV capacity in blocks.
\end{itemize}
The learned policy reads nothing else. In particular it ignores the request's expected output
length, which offline replay fills with the true output length, and every performance-model
estimate; engine-internal state such as waiting-queue length or KV usage is not available to it.
Because it uses only these inputs, the artifact trained in replay is the one the live router runs
(\cref{sec:integration}). Replay applies cache events to the
router's index at once, while a live router sees them late; this difference is a robustness check
(\cref{sec:robustness}), not part of the problem definition.

\subsection{Load modes}
\label{sec:loadmodes}

% src: CR/cells/{train,val,test}.jsonl load.mode by family (train: Mooncake 12 open + 6 closed,
%      sessions 6 open + 2 closed, AgentX 8 lanes; test: 31 open, 15 closed, 14 lanes);
%      CR/facts/calibration.json rules.levels (units), lessons.LR-09 (causal session release kept);
%      WT/benchmarks/learned_routing/learned_routing/goodput.py (occupancy units); CR/CONTRACT.md A3
A cell offers its trace to the workers in one of three ways, which differ in what a policy can
change.
\begin{description}
  \item[Open loop.] Requests arrive at the trace's timestamps, compressed or stretched by a speedup
    factor that sets the load. A policy cannot change when requests arrive, only how well they are
    served. In multi-turn session cells the session start times are fixed this way, while each
    follow-up turn is released a think time after the previous turn completes, as in a real
    conversation.
  \item[Closed loop.] A fixed number of sessions is in flight; a session holds its slot through its
    think times, and a new session starts when one ends. Each Mooncake trace row is its own
    one-turn session. A policy that serves requests faster also receives more of them, so
    throughput depends on the policy.
  \item[Agentic lanes.] A fixed number of lanes each run one recorded agent play at a time. Inside
    a play, requests follow the recorded dependency graph (sequential turns, subagent spawns and
    joins, tool delays), and a lane starts its next play when the previous one ends.
\end{description}
Mooncake, FAST25 and session cells run in open and closed loop, and AgentX cells in lanes;
calibration sets each cell's load (\cref{sec:calibration}). Results are reported per mode, because
open and closed systems respond to scheduling differently \citep{schroeder2006open}. Session and
agentic work is released causally because replaying follow-up turns at their recorded timestamps
would fix a release schedule, and cache state, that a different policy is supposed to change
\citep{rajib2026agentservesim,alomar2023causalsim}.
% src: CR/literature/LESSONS.md LR-09 (AgentServeSim pdf p.1-2, p.4; CausalSim pdf p.1-3;
%      Schroeder pdf p.10-11), VERIFIED

\subsection{Objective}
\label{sec:objective}

\paragraph{Good requests.}
% src: CR/CONTRACT.md A2.2 (operator decision: no TTFT SLO); WT/benchmarks/learned_routing/
%      learned_routing/goodput.py (a2_good, token_form_good, completed, E2E_REL_TOL = 1e-6)
The deployment's service objective, set by the operator, concerns the pace of the token stream and
the total request latency, not the time to first token. A request is \emph{good} iff it completes
and
\begin{equation}
  \label{eq:good}
  \good_j = \mathbf{1}\!\left[\ITLbar_j \le \Imax\right]\cdot
            \mathbf{1}\!\left[\EtoE_j \le \Smax\,\Ezero(\ISL_j,\OSL_j)\right],
\end{equation}
where $\ITLbar_j = (\EtoE_j - \TTFT_j)/(\OSL_j - 1)$ is the mean inter-token latency, the ITL
clause is skipped when $\OSL_j \le 1$, and the slowdown clause allows a relative tolerance of
$10^{-6}$ for floating-point rounding. Rejected, errored and incomplete requests are never good.
$(\Imax, \Smax)$ is one pair per workload family (\cref{tab:notation}). $\TTFT_j$ is reported but
is not part of \cref{eq:good}; it enters only through $\EtoE_j$.

\paragraph{The reference latency.}
% src: WT/benchmarks/learned_routing/learned_routing/e0.py (method ais-chunked-estimator-v2;
%      "matches 17,233 isolated single-request replays to within 2e-8 relative");
%      CR/facts/STATE.md audit(calibration/independent-rerun r1) (own E0 == harness within
%      1.11e-8 on 880,139 pairs); CR/CONTRACT.md A2.2
$\Ezero(\ISL,\OSL)$ is the request's latency alone on an idle worker with no prefix reuse: a full
chunked prefill of $\ISL$ tokens followed by $\OSL - 1$ decode steps at batch size 1, timed by
\ais{}. It depends only on $(\ISL, \OSL)$ and is cached per pair. It agrees with isolated
single-request replays to within $2\times10^{-8}$ relative, and an independent audit's own
computation matched it on \num{880139} pairs. $\Ezero$ is used only for scoring; no policy sees it.

\paragraph{Why a slowdown bound.}
% src: CR/facts/setup.json long_context.evidence.single_request_ttft_ms_ais (1024: 55.49 ms,
%      8192: 475.97 ms, 32768: 2,550.45 ms, 120000: 17,815.54 ms);
%      CR/facts/build_workloads.json traces.flat.<family>.isl_mean (mooncake 8,590.0;
%      fast25_conversation 12,035.1; fast25_synthetic 15,325.5); agentx.selection.predicate
%      (in + max(out, 1) <= 131072); CR/literature/LESSONS.md LR-08 (SMetric pdf p.8, VERIFIED;
%      CR/literature/txt/llm-routing/wang2026-smetric.txt line 1135)
Request sizes vary widely. Mean prompt lengths are \num{8590} tokens on Mooncake
and \num{12035} to \num{15326} on the FAST25 traces, and AgentX plays were admitted as long as
every request, prompt plus output, fits the \num{131072}-token context. In \ais{}'s timing, prefill
alone on an idle worker takes \SI{55.49}{ms} at \num{1024} tokens, \SI{475.97}{ms} at \num{8192},
\SI{2550}{ms} at \num{32768} and \SI{17815.54}{ms} at \num{120000}. A single absolute bound on
end-to-end latency would be unattainable for the longest requests even on an idle worker, or slack
for the shortest. Fixed first-token deadlines have the same problem; SMetric's authors note that
providers typically set the TTFT deadline linear in the request length \citep{wang2026smetric}. Dividing by $\Ezero$ scales the bound
with each request's own unavoidable work: under the Mooncake bound $\Smax = 3.46625$, a request may
take up to about 3.5 times its uncontended latency, whatever its size.

Two properties of $\Ezero$ keep the bound fair across policies. It is policy independent, so every
policy is held to the same per-request bound. And it assumes no prefix reuse, so a cache hit lowers
$\EtoE_j$ but not $\Ezero$: a policy that earns more reuse finds the bound easier to meet, never
harder.

\paragraph{Why two clauses.}
The clauses catch different failures. For a long prompt, $\Smax\Ezero$ is dominated by prefill
time, so a slow token stream after a fast prefill can pass the slowdown clause; the ITL clause
catches it. Conversely, a request that waits a long time and then streams at an ordinary pace
passes the ITL clause and fails the slowdown clause. At the calibrated bounds the slowdown clause
does most of the work: at the calibration load nearest L2 on $\Nw = 8$, the ITL clause alone fails
0.6\% of $\piref$'s requests on sessions, 1.8\% on AgentX and 4.5\% on Mooncake, while 15\% of
requests are not good at L2 by construction.
% src: CR/facts/calibration.json sla.<family>.itl_only_fail_at_l2_grid.default
%      (synthetic_sessions 0.00625 at grid load 0.375; agentx 0.0178 at 3; mooncake 0.0449 at 9000)
%      and summary.sla.<family>.itl_only_fail_frac_default_at_L2; L2 = good fraction 0.85
%      (rules.levels). "Most of the work" is the inference 15% - 4.5% > 4.5%.
% check: the grid load is near but not at L2 (sessions 0.375 vs 0.380, AgentX 3 vs 3.388,
%        Mooncake 9000 vs 8797.97 per worker); keep "nearest L2".

\paragraph{Windowed goodput.}
% src: CR/CONTRACT.md A2.3; CR/facts/calibration.json rules.windows;
%      WT/benchmarks/learned_routing/learned_routing/goodput.py (measurement_window, _window_counts,
%      _completes_inside, _completion_end)
A cell's score is its windowed goodput $\Gw_{\cell,\krep}(\pol)$, good requests per simulated
second inside a measurement window $\win_{\cell}$ that starts after a warm-up. Its form depends on
the load mode. In open-loop cells the window is a fixed interval of arrival times, and
\begin{equation}
  \label{eq:goodput-open}
  \Gw_{\cell,\krep}(\pol) = \frac{1}{|\win_{\cell}|}\sum_{j:\ a_j \in \win_{\cell}} \good_j ,
\end{equation}
so a request that arrives in the window counts even if it finishes after it, and one that never
completes counts as not good. In closed-loop and lane cells, $\Gw$ counts the good requests that
\emph{complete} inside a half-open window, divided by its length. That window ends at full
occupancy, the last instant at which every session or lane slot is busy, so its end, and hence
its length, depends on the policy.

The window removes three distortions.
\begin{itemize}
  \item \emph{Drain.} The replay's native goodput divides good requests by the makespan, which
    includes the drain after the last arrival, so it mixes SLO attainment with how fast a policy
    clears the tail.
    % src: CR/literature/LESSONS.md LR-01 [code] aisimulate-core src/replay/report.rs 1807-1895
    %      (duration_ms = max(report_end, max terminal_time) - report_start), VERIFIED
    In closed loops and lanes, completions after full occupancy happen with fewer requests in
    flight, and in lane cells the lanes run out of plays at different times, so counting the tail
    would credit easy, under-loaded requests.
    % src: goodput.py module docstring "Completion-basis end rule (build audit r1 F1)"
  \item \emph{Cold start.} Caches start empty and session populations ramp up, so the first minutes
    of a replay are not the steady state the cell represents. The warm-up length matters: for
    open-loop session cells, a first calibration with a \SI{180}{s} warm-up gave $\piref$ a good
    fraction that moved by up to 0.353 when the warm-up history was doubled (9 of 12 cells beyond
    3 standard errors); with the adopted warm-up of two session lifetimes, doubling moves it by at
    most 0.014 (no cell beyond 3 standard errors).
    % src: CR/facts/calibration.json session_open_warmup.doubled_warmup_check.summary
    %      ("w180_vs_w2x|gf_default": max_abs_delta 0.3527, cells_abs_z_gt_3 9 of 12;
    %       "w1x_vs_w2x|gf_default": max_abs_delta 0.0142, cells_abs_z_gt_3 0 of 12)
  \item \emph{Deferral.} An objective that ignored requests finishing after the horizon would
    reward deferring work past it; Decima's learned scheduler did exactly this under a fixed
    horizon \citep{mao2019decima}. Counting in-window arrivals that finish late removes the
    incentive.
    % src: CR/literature/LESSONS.md LR-01 (Decima pdf p.7), VERIFIED
\end{itemize}
When every arrival is set by the trace, as in the open-loop Mooncake and FAST25 cells, the number
of in-window arrivals is the same for every policy on a replicate, so the ratio of two policies'
goodputs equals the ratio of their in-window good fractions. In closed-loop and lane cells
goodput also rewards throughput, which the policy legitimately controls there. SLO attainment as
the primary serving metric follows \citet{zhong2024distserve}.

\paragraph{SLO and load calibration.}
% src: CR/facts/calibration.json rules.sla, rules.levels, sla.<family>; CR/runs/calibrate-fix-r0/decisions.json
The bounds $(\Imax,\Smax)$ are fixed per family from runs of $\piref$ on train segments at
$\Nw = 8$, and each cell's load is one of three levels L1, L2, L3, at which $\piref$'s good fraction
is 0.95, 0.85 and 0.65. Levels are calibrated separately for every $\Nw$ and transform, on train
segments only, and both bounds and levels were frozen before any policy was tuned
(\cref{sec:calibration}). The bounds differ widely by family, from $\Smax = 1.16134$ on AgentX to
3.46625 on Mooncake, because each is anchored at its own family's load knee. Every cell therefore
sits where the default router is starting to miss its objective, the regime in which routing
changes outcomes; that the regime is defined relative to one baseline is a threat discussed in
\cref{sec:threats}.

\paragraph{Normalized score.}
% src: WT/benchmarks/learned_routing/learned_routing/train.py (pair_score, LOG_CLIP = ln 3,
%      TrainConfig.eps = 1e-3); CR/facts/HEADLINE_TEST.json per_cell_statistic;
%      CR/facts/gate.json full_plan.objective_cells_repeats ("pre-registered, LR-01")
Goodput levels differ by orders of magnitude across cells, so every comparison is a ratio on the
same cell and replicate:
\begin{equation}
  \label{eq:clr}
  \ratio_{\cell,\krep}(\pol,\pol') = \frac{\Gw_{\cell,\krep}(\pol)+\varepsilon}{\Gw_{\cell,\krep}(\pol')+\varepsilon},
  \qquad
  \clr_{\cell,\krep}(\pol,\pol') = \clip\!\left(\log \ratio_{\cell,\krep}(\pol,\pol'),\,-\clipb,\,\clipb\right),
\end{equation}
with $\varepsilon = 10^{-3}$ good requests per second. Both policies run on the same workload
replicate $\krep$ with the same tie-breaking seed (common random numbers, \cref{sec:noise}).
Training uses $\pol' = \piref$ (\cref{eq:objective}); the headline uses $\pol' = \pibase$
(\cref{eq:segstat}). The log makes the score symmetric, so doubling goodput on one cell and halving
it on another nets to zero. The clip limits any one cell to a factor of 3 either way, which matters
because a mean of raw ratios is dominated by the few cells where one policy collapses
\citep{agarwal2021precipice}, and such cells occur here. On all 9 open-loop session cells of the
train and validation splits, round-robin's goodput is 0.111 to 0.425 times the default router's,
and below 1/3 on 5 of them. A policy that far past its stability limit has no steady state, so the
ratio is not a stable quantity: on one cell, doubling the warm-up history halved it, from 0.118 to
0.055. The clip scores every ratio below 1/3 alike, so that unstable detail cannot drive training
or the headline.
% src: CR/audits/calibration/independent-rerun-r1.md F1 ("RR/default ratio is 0.111-0.425; 5 train/val
%      cells are below 1/3"; sessions-s1-base-n4-open-L3 0.118 -> 0.055 with doubled history;
%      "absorbs everything below 1/3"); Agarwal pdf p.2-3 (CR/literature/LESSONS.md LR-01, VERIFIED)
The clipped log-ratio was chosen over the mean of ratios at calibration, before any training.
% src: CR/facts/HEADLINE_TEST.json (written "calibration stage, before any training");
%      CR/facts/pilot.json decisions[1] (objective clipped_log_ratio, "pre-registered in HEADLINE_TEST.json")

\paragraph{What ``better'' means.}
% src: CR/CONTRACT.md A2.1, A17.2; CR/facts/HEADLINE_TEST.json; CR/CONTRACT.md A14
The headline compares the learned policy with the best baseline \emph{on the same footing}: each
policy, learned or heuristic, is one configuration tuned on the pooled train split with the same
CMA-ES budget and selected on validation (\cref{sec:training,sec:eval}). The best baseline $\pibase$
is the baseline with the highest mean clipped log-ratio on fresh validation replicates. The learned
side is chosen by the same rule among the learned variants that read only router-observable
signals, and the number of variants compared is reported. The comparison is made once, on the
frozen test set, by the pre-registered test of \cref{sec:headline}. Because the objective itself is
a modeling choice, the same test records are also re-scored with scaled bounds and with
alternative definitions of a good request: ITL alone, slowdown alone, and the ITL bound combined
with either a length-scaled TTFT bound or an absolute latency bound. These secondary analyses never
change the headline.

% SPDX-FileCopyrightText: Copyright (c) 2026 NVIDIA CORPORATION & AFFILIATES. All rights reserved.
% SPDX-License-Identifier: Apache-2.0
%
% The single notation table. Every section uses these symbols through the macros in
% macros.tex. A new symbol needs a row here first.
% Code names refer to WT/lib/router-plugins/builtin/src/learned_choice/ (Rust) and
% WT/benchmarks/learned_routing/learned_routing/ (harness); contract names to CR/CONTRACT.md.

\subsection{Notation}
\label{sec:notation}

\Cref{tab:notation} fixes the symbols used throughout. The campaign's contract, the policy's
specification (\code{FEATURES.md}) and the harness each grew their own names, and several
collide. We resolve the collisions as follows.

\begin{itemize}
  \item \textbf{The letter $S$.} $\cand$ is the set of eligible candidate workers for one request,
    as in \code{FEATURES.md} ($\fbar$, $\ctxv_{\cand,\cidx}$). The contract's end-to-end
    slowdown bound, also called $S$ there, is written $\Smax$ here. For symmetry the contract's
    mean-ITL bound $I$ is written $\Imax$.
  \item \textbf{The letter $k$.} $\krep$ indexes common-random-number (CRN) workload replicates,
    as in the contract (Amendment A1) and the harness (\code{repeat}). \code{FEATURES.md} indexes
    context terms by $k$; here they are indexed by $\cidx = 1,\dots,\rank$.
  \item \textbf{The letter $B$.} $\budget$ is the CMA-ES evaluation budget per run. The
    request's block count, which \code{FEATURES.md} calls $B$, is $\reqblocks_j$ here.
  \item \textbf{The letter $T$.} \code{FEATURES.md}'s token normalizer $T = 8192$ is $\Tch$.
    The plan's TTFT threshold $T$ no longer exists (Amendment A2 dropped the TTFT SLO).
  \item \textbf{$e$ and $E$.} $\basis{f}$ (bold) is a standard basis vector, so the parity anchor
    is $\theta = -\basis{0}$; $\Ezero$ is a latency.
  \item \textbf{$\tau$.} $\temp$ is the policy temperature; Kendall's rank correlation is
    $\kendall$.
  \item \textbf{$\sigma$ and segments.} $\sigma$ belongs to CMA-ES (step size $\sigma$, initial
    $\sigmazero$); workload segments are $\seg \in \Omega$. Ratios are $\ratio$, restarts $\restarts$.
\end{itemize}

{\small
\setlength{\LTcapwidth}{\linewidth}
\begin{longtable}{@{}P{0.17\linewidth}P{0.44\linewidth}P{0.33\linewidth}@{}}
\caption{Notation. ``Value'' is the setting in this campaign; code and file names are in
  typewriter. Pilot values are \prelim.}
\label{tab:notation}\\
\toprule
Symbol & Meaning & Value / code name \\
\midrule
\endfirsthead
\toprule
Symbol & Meaning & Value / code name \\
\midrule
\endhead
\bottomrule
\endfoot
\multicolumn{3}{@{}l}{\emph{Deployment and requests}}\\
% src: CR/CONTRACT.md "Deployment (operator-fixed)"; CR/cells/SPLIT_MANIFEST.json
$\Nw$ & number of identical workers (one Qwen3-32B TP2 vLLM engine each) &
  $\{2,4,6,8,16,32\}$; train $\{4,8\}$, val $\{4,6,8\}$, test all; \code{num\_workers} \\
$j$ & request index & \\
$a_j$ & arrival time of request $j$ (simulated) & \\
$\ISL_j$, $\OSL_j$ & prompt and output length (tokens) &
  \code{input\_length}, \code{output\_length} \\
% src: WT/lib/router-plugins/builtin/src/learned_choice/FEATURES.md "Feature set v1"
$\reqblocks_j$ & request blocks, $\max(\lceil \ISL_j / 16\rceil, 1)$ (\code{FEATURES.md}: $B$) &
  KV block size 16; \code{request\_blocks} \\
$\TTFT_j$, $\EtoE_j$ & time to first token; end-to-end latency (simulated) & \code{per\_request} rows \\
$\ITLbar_j$ & mean inter-token latency, $(\EtoE_j - \TTFT_j)/(\OSL_j - 1)$, for $\OSL_j > 1$ & \\
% src: CR/CONTRACT.md A2.2; CR/facts/calibration.json (E0 method)
$\Ezero(\ISL,\OSL)$ & uncontended no-reuse latency: the request alone on an idle worker, full
  prefill, then $\OSL$ decode steps at batch 1, timed by \ais. Evaluation only; no policy sees it &
  \code{ais-chunked-estimator-v2}; \code{learned\_routing/e0.py} \\
% src: CR/facts/calibration.json sla.<family>.itl_ms
$\Imax$ & mean-ITL bound of the SLO (contract: $I$) &
  Mooncake and FAST25 \SI{252.117}{ms}; sessions \SI{46.0478}{ms}; AgentX \SI{38.6205}{ms} \\
% src: CR/facts/calibration.json sla.<family>.e2e_slowdown
$\Smax$ & end-to-end slowdown bound (contract: $S$) &
  Mooncake and FAST25 3.46625; sessions 2.11784; AgentX 1.16134 \\
$\good_j$ & good-request indicator, \cref{eq:good} & \code{learned\_routing/goodput.py} \\
\midrule
\multicolumn{3}{@{}l}{\emph{Policy}}\\
$\cand$ & eligible candidate set of one request; $|\cand| \le \Nw$ & host eligibility \\
$i$ & a candidate worker, $i \in \cand$ & \\
$\fvec_i \in \mathbb{R}^{\dfeat}$ & feature vector of candidate $i$, fixed order (\cref{tab:features}) &
  $\dfeat = 8$ (\code{v1}) or 23 (\code{v2}); \code{feature\_set} \\
$\fbar$ & candidate mean of $\fvec$ over $\cand$ & \\
$\basis{f}$ & standard basis vector; $\theta = -\basis{0}$ reproduces the default router & \\
$\Tch$ & token normalizer of the prefill features (\code{FEATURES.md}: $T$) & 8192 \\
$\theta \in \mathbb{R}^{\dfeat}$ & linear coefficients & \code{parameters.theta}; $\theta_0 = -1$ pinned in training \\
$\rank$, $\cidx$ & context rank; context-term index $\cidx = 1,\dots,\rank$ (\code{FEATURES.md}: $k$) &
  $\rank = 0$ (M1), $\rank = 2$ (M2); \code{len(context.p)} \\
$p_{\cidx} \in \mathbb{R}^{\dfeat}$ & context direction of term $\cidx$ & \code{context.p} \\
$\ctxv_{\cand,\cidx}$ & context value of term $\cidx$: a named source (\cref{tab:sources}) or, in the
  pooled form, $q_{\cidx}^{\top}\fbar$ & \code{context.sources} or \code{context.q} \\
$\weff$ & effective weights, $\theta + \sum_{\cidx} \ctxv_{\cand,\cidx}\, p_{\cidx}$ & \code{features.rs} \\
$\util_i$ & utility of candidate $i$, $\weff^{\top}\fvec_i$ & \\
$\temp$ & temperature: $\temp = 0$ is argmax with a seeded tie-break, $\temp > 0$ samples
  $\softmax(\util/\temp)$ & \code{temperature}; fixed at 0 for M1 and M2 \\
$\pol$, $\pilearn$ & a routing policy; the learned policy & one policy YAML each \\
$\piref$ & \defaultref: \defaultcf{} at shipped parameters; the normalizer of every ratio &
  harness name \code{default} \\
$\pibase$ & the validation-selected best baseline (\cref{sec:eval}) & tuned \policy{ramjet} \\
\midrule
\multicolumn{3}{@{}l}{\emph{Workloads, metric and tests}}\\
$\cell$ & cell: one trace segment, transform, load mode and value, $\Nw$ and SLO & \code{cell\_id} \\
% src: CR/cells/{train,val,test}.jsonl line counts
$\Cset{train}, \Cset{val}, \Cset{test}$ & cells of each split & 34, 14, 60 \\
% src: CR/cells/SPLIT_MANIFEST.json segments_by_split
$\seg$, $\Segset{split}$ & segment (unit of statistical independence); segments of a split &
  10, 6, 12 \\
$\win_{\cell}$ & measurement window of cell $\cell$ (\cref{sec:objective}) & \code{measure} \\
L1, L2, L3 & load levels at which $\piref$'s good fraction is 0.95, 0.85, 0.65 & \code{load.level} \\
$\krep$ & CRN replicate index; replicate $\krep$ also sets policy seed $\krep + 1$ & \code{repeat} \\
$\Gw_{\cell,\krep}(\pol)$ & windowed goodput (good requests per simulated second) &
  \code{goodput\_rps\_window} \\
$\varepsilon$ & offset in goodput ratios & $10^{-3}$; \code{--eps} \\
$\ratio_{\cell,\krep}(\pol,\pol')$ & $(\Gw_{\cell,\krep}(\pol)+\varepsilon)/(\Gw_{\cell,\krep}(\pol')+\varepsilon)$;
  $\pol' = \piref$ unless stated & \code{pair\_score} \\
$\clr_{\cell,\krep}(\pol,\pol')$ & clipped log-ratio, $\clip(\log\ratio, -\clipb, \clipb)$ &
  \code{clipped\_log\_ratio} \\
$\segstat_{\seg}$ & per-segment headline statistic, \cref{eq:segstat} &
  \code{facts/HEADLINE\_TEST.json} \\
% src: CR/facts/pilot.json mde.learned_vs_val_best_default_arm_heuristic.mde (0.03815); CR/facts/HEADLINE_TEST.json
MDE & minimum detectable effect, $\max(1\%,\ 2 \times$ validation segment SE$)$ &
  0.038, preliminary (pilot, validation split) \\
% src: CR/CONTRACT.md A5; CR/facts/lag_patch.json
$\lag$ & router-state lag, robustness check only & $\{10, 50, 200\}$\,ms on test; \code{router\_state\_lag\_ms} \\
$\kendall$ & Kendall rank correlation of two policy rankings (a perturbed against the nominal
  condition; simulation against live GPUs) & $\tau$-b \\
\midrule
\multicolumn{3}{@{}l}{\emph{Training}}\\
$\Jobj(\pol)$ & training objective, \cref{eq:objective} & \code{--objective clipped\_log\_ratio} \\
$\searchvec$ & CMA-ES internal coordinates of the free parameters (bounded ones in $[0,1]$) &
  \code{space.yaml} \\
$\gen$ & CMA-ES generation & $0,\dots,24$ \\
% src: CR/facts/gate.json full_plan.budget
$\popsize$ & population size & 16; \code{--popsize} \\
$\budget$ & objective evaluations per run (\code{FEATURES.md}'s $B$ is $\reqblocks_j$) &
  $400 = 25 \times 16$; \code{--budget-evals} \\
$\Krep$, $\Kpool$ & replicates per objective evaluation; the pool they rotate through &
  2 and 8; \code{--replicates}, \code{--replicate-pool} \\
$\Kset{\gen}$ & replicates of generation $\gen$, $\{(\gen\Krep + i) \bmod \Kpool : i < \Krep\}$ &
  \code{Trainer.eval\_set} \\
% src: CR/facts/gate.json full_plan.objective_cells_repeats; train.py TrainConfig
$\Kval$ & replicates per validation score inside \code{lr-train} & 3 ($\krep = 0,1,2$) \\*
$\sigmazero$ & initial CMA-ES step size, in internal units & per policy (\cref{sec:training}) \\*
$\restarts$ & restarts per policy & 3 \\
\end{longtable}
}

